\documentclass[11pt]{article}
\usepackage[T1]{fontenc}
\usepackage{textcomp}
\usepackage[margin=0.75in]{geometry}
\usepackage{amsmath,amssymb,amsthm}
\usepackage{array,booktabs}
\usepackage{hyperref}
\setlength{\parindent}{0pt}
\newtheorem{theorem}{Theorem}
\theoremstyle{definition}
\newtheorem{definition}{Definition}
\title{Flow Proof Artifact: List-len-cons-cons}
\author{Generated by \texttt{flow doc proof}}
\date{Flow Proof Book}
\begin{document}
\maketitle
A two-element list has length two.\\[0.5em]
\textbf{Source.} church — https://en.wikipedia.org/wiki/Length\_of\_a\_list\\[0.5em]
\bigskip
\noindent{\large \textbf{Derived fact 1.} \textit{len(cons(x, cons(y, nil))) = succ(succ(0))} \hfill List \cdot length \cdot pair list has length two}\par
\smallskip\noindent\textit{Coordinate: List \cdot length \cdot pair list has length two}\par
\smallskip\noindent\textit{Source: church}\par
\smallskip\noindent\textit{Built on: cons increases length by one, for length on List, singleton has length one, for length on List}\par
\medskip
\noindent\textbf{Goal.} len(cons(x, cons(y, nil))) = succ(succ(0))\par
\noindent$\displaystyle \forall x \in \mathbb{N} \forall y \in \mathbb{N}\quad len(cons(x, cons(y, nil))) = \mathrm{succ}(succ(0))$\par
\medskip
\noindent
\begin{tabular}{@{} >{\raggedright\arraybackslash}p{0.06\textwidth} >{\raggedright\arraybackslash}p{0.47\textwidth} >{\raggedleft\arraybackslash}p{0.41\textwidth} @{}}
\textbf{\#} & \textbf{Proof} & \textbf{Mathematics} \\
\midrule
\textbf{1.} & We prove that pair list has length two for length on List. & \\[0.45em]
\textbf{2.} & We invoke the derived fact governing length on List: cons increases length by one, for length on List (instantiated for x, cons(y, nil)). & \\[0.45em]
\textbf{3.} & We invoke the derived fact governing length on List: singleton has length one, for length on List (instantiated for y). & \\[0.45em]
\textbf{4.} & From step 2 and step 3, this implies len(cons(x, cons(y, nil))) equals the successor of succ(0). Hence proven. & $\displaystyle len(cons(x, cons(y, nil))) = \mathrm{succ}(succ(0))$ \\[0.45em]
\bottomrule
\end{tabular}
\label{thm:List-length-pair_list_has_length_two}
\par\medskip\hrule
\end{document}
